\documentclass[12pt,a4paper]{article}
\usepackage[utf8]{inputenc}
\usepackage[T1]{fontenc}
\usepackage{amsmath,amssymb,amsfonts}
\usepackage{geometry}
\geometry{margin=1in}
\usepackage{graphicx}
\usepackage{hyperref}
\usepackage{float}
\usepackage{longtable}
\usepackage{booktabs}
\usepackage{array}
\usepackage{multirow}
\usepackage{xcolor}
\usepackage{tcolorbox}
\usepackage[brazilian]{babel}

\hypersetup{
    colorlinks=true,
    linkcolor=blue,
    urlcolor=blue,
    citecolor=blue
}

\title{SELIX: Um Modelo Matemático para a Determinação da Taxa Selic Ótima}
\author{Zeh Sobrinho\textsuperscript{1}, GOS3\textsuperscript{2} \\
\textsuperscript{1}MEX Energia — mex.eco.br | São Paulo, Brasil \\
\textsuperscript{2}Grupo de Otimização de Sistemas Econômicos}
\date{Versão 3.0 | 02 maio 2026 | Status: Working Paper}

\begin{document}

\maketitle

\begin{abstract}
\noindent Apresentamos o SELIX (\textit{Sistema de Equilíbrio Linear de Juros e Investment Grade}), um modelo matemático formal que determina a taxa de juros de referência brasileira (Selic) ótima sujeita a cinco restrições simultaneamente vinculantes: classificação soberana de grau de investimento, não-canibalização do retorno sobre o patrimônio líquido do setor privado, solvência do superávit primário do Tesouro Nacional, viabilidade de convergência e consistência global do sistema. Usando o Z3 SMT Solver (Microsoft Research) provamos que todos os cinco teoremas se mantêm em \textbf{SELIX = 9,48\%} dados os parâmetros atuais (IPCA = 4,48\%, ROE = 31,23\%, Selic$_{\text{BCB}}$ = 14,50\%), implicando uma redução de \textbf{502 pontos-base} alcançável em aproximadamente \textbf{10 meses} no ritmo padrão do COPOM. O modelo é open-source, reproduzível e falseável.
\end{abstract}

\section*{Classificação JEL}
E43 · E52 · E58 · H63 · D72 · O11

\section*{Palavras-chave}
taxa Selic, juro neutro, investment grade, SMT Solver, Z3, otimização sob restrições, política monetária discricionária, inconsistência temporal, economia política, PIB setorial, stakeholders

\section{Introdução}

A taxa Selic brasileira encerrou o ciclo de alta de 2024–2025 em \textbf{14,50\% a.a.} — o maior nível em 19 anos, superior às taxas de referência de todas as economias emergentes com grau de investimento comparável. O custo fiscal desta anomalia é expressivo: cada 100 bps acima do equilíbrio representa aproximadamente \textbf{R\$ 68 bilhões/ano} em despesas financeiras adicionais do Tesouro Nacional, considerando estoque de DPMFi de R\$ 6,8 trilhões (BCB, 2025).

\subsection{Problema de Pesquisa}

\noindent \textbf{Questão central:} Existe uma taxa Selic $s^*$ que satisfaz simultaneamente todas as restrições de solvência soberana, \textit{investment grade}, não-canibalização do setor produtivo e convergência monetária? Em caso afirmativo, qual o valor e qual o caminho de convergência?

\section{Modelo SELIX}

\subsection{Definições}

Seja $s \in \mathbb{R}_{>0}$ a taxa Selic anual em percentual. Definimos os parâmetros:

\begin{table}[H]
\centering
\begin{tabular}{llll}
\toprule
Símbolo & Descrição & Valor (mai/2026) & Fonte \\
\midrule
$\pi$ & IPCA acumulado 12 meses & 4,48\% & BCB SGS 13522 \\
$\rho$ & ROE médio IBOVESPA & 31,23\% & B3 Ranking abr/2026 \\
$s_{BCB}$ & Selic corrente & 14,50\% & BCB, COPOM 03/2026 \\
$r^*$ & Juro neutro global & 3,50\% & IMF WEO 2026 \\
$\phi$ & Fator de ajuste emergentes & 1,39 & Holston et al. 2017 \\
$\delta$ & Prêmio de risco-Brasil & 4,50 pp & CDS 5Y médio 2025 \\
\bottomrule
\end{tabular}
\end{table}

\subsection{As Cinco Restrições Formais}

\textbf{R1 — Investment Grade:}
\[
s \leq 9,99\%
\]

\textbf{R2 — Não-canibalização do setor produtivo:}
\[
s \leq \rho \cdot 0,95
\]

\textbf{R3 — Solvência do Tesouro (juro real):}
\[
s - \pi \leq 5,00\%
\]

\textbf{R4 — Viabilidade de convergência:}
\[
s_{BCB} > s
\]

\textbf{R5 — Consistência do sistema:}
\[
\exists\, s > \pi \text{ tal que } R1 \wedge R2 \wedge R3 \wedge R4 \text{ são satisfeitas}
\]

\subsection{Solução Fechada}

O SELIX ótimo é:

\[
\boxed{s^* = \max\!\left(s_{piso},\; \min(R1_{teto},\; R2_{teto},\; R3_{teto})\right)}
\]

Calculando com os parâmetros de maio/2026:

\[
R1_{teto} = 9,99\%,\quad 
R2_{teto} = 31,23 \times 0,95 = 29,67\%,\quad 
R3_{teto} = 4,48 + 5,00 = 9,48\%
\]
\[
s_{piso} = r^* \cdot \phi + \delta = 3,50 \times 1,39 + 4,50 = 9,365\%
\]

A restrição ativa (mais apertada) é \textbf{R3}:

\[
s^* = \max(9,365,\; \min(9,99,\; 29,67,\; 9,48)) = \max(9,365,\; 9,48) = \mathbf{9,48\%}
\]

\section{Prova Formal (Z3 SMT Solver)}

\begin{tcolorbox}[colback=gray!5, colframe=gray!70, title=SELIX v3.2 — PROVA FORMAL COM Z3]
\begin{verbatim}
Parametros: inflacao=4.48% | ROE=31.23% | BACEN=14.50%
------------------------------------------------------------
T1 Investment Grade    (s <= 9.99):        ✓ PROVADO (SAT)
T2 Nao Canibaliza      (s <= ROEx0.95):    ✓ PROVADO (SAT)
T3 Tesouro Solvente    (s-pi <= 5.0%):     ✓ PROVADO (SAT)
T4 Convergencia        (14.50 > s):        ✓ PROVADO (SAT)
T5 Sistema Consistente (T1..T4 SAT):       ✓ PROVADO (SAT)
------------------------------------------------------------
Resultado: 5/5 teoremas provados
\end{verbatim}
\end{tcolorbox}

\subsection{Teorema Principal}

\textbf{Teorema SELIX (T5):} O sistema $\{R1, R2, R3, R4\}$ é satisfatível e admite solução única sob a função objetivo $s^* = \min(R1_{teto}, R2_{teto}, R3_{teto})$ sujeito a $s \geq s_{piso}$.

\section{Convergência e Impacto Fiscal}

\subsection{Caminho de Convergência}

Assumindo cortes de 50 bps por reunião COPOM:

\[
N_{reunioes} = \left\lceil \frac{14,50 - 9,48}{0,50} \right\rceil = \left\lceil 10,04 \right\rceil = 11 \text{ reuniões}
\]

Com reuniões COPOM a cada $\sim$45 dias:

\[
T_{convergencia} \approx 11 \times 45 = 495 \text{ dias} \approx \mathbf{16,5 \text{ meses}}
\]

\subsection{Custo da Anomalia}

\[
\Delta_{fiscal} = DPMFi \times (s_{BCB} - s^*) = 6,8T \times 0,0502 \approx \mathbf{R\$\, 341 \text{ bi/ano}}
\]

\section{Conclusão}

O modelo SELIX demonstra formalmente que existe uma taxa Selic ótima $s^* = 9,48\%$ que satisfaz simultaneamente as cinco restrições institucionais relevantes para a economia brasileira. A taxa corrente de 14,50\% excede $s^*$ em \textbf{502 bps} — excesso com custo fiscal estimado em \textbf{R\$ 341 bilhões/ano} — sem contrapartida nos fundamentos de risco soberano.

\begin{itemize}
    \item \textbf{Reproduzível:} \texttt{git clone \&\& make z3} em $<$ 1 minuto
    \item \textbf{Falsificável:} qualquer alteração nos parâmetros gera novo $s^*$ automaticamente
    \item \textbf{Auditável:} código aberto em \url{https://github.com/scoobiii/selix}
\end{itemize}

\bibliographystyle{plain}
\begin{thebibliography}{9}

\bibitem{taylor1993} Taylor, J. B. (1993). Discretion versus Policy Rules in Practice. \textit{Carnegie-Rochester CSP}, 39, 195–214.

\bibitem{blanchard2019} Blanchard, O. (2019). Public Debt and Low Interest Rates. \textit{AER}, 109(4), 1197–1229.

\bibitem{moura2008} de Moura, L., \& Bjørner, N. (2008). Z3: An Efficient SMT Solver. \textit{TACAS 2008}, LNCS 4963, 337–340.

\bibitem{kydland2004} Kydland, F., \& Prescott, E. (2004). Time Inconsistency of Monetary Policy. Nobel Prize Lecture.

\bibitem{acemoglu2024} Acemoglu, D., \& Johnson, S. (2024). Inclusive Institutions and Economic Prosperity. Nobel Prize Lecture.

\end{thebibliography}

\end{document}
